\inputencoding{utf8}
\HeaderA{loca\_t\_rast}{Load LOCA2 tmax and tmin files into rasters}{loca.Rul.t.Rul.rast}
%
\begin{Description}
Load LOCA2 tmax and tmin files into rasters
\end{Description}
%
\begin{Usage}
\begin{verbatim}
loca_t_rast(filepath, scenario = c("historical", "ssp585"))
\end{verbatim}
\end{Usage}
%
\begin{Arguments}
\begin{ldescription}
\item[\code{filepath}] file path where files are located. This should
be the folder in which the climate scenario folders are contained
(historical, ssp585, etc.). For example, "ACCESS-CM2/0p0625deg/r1i1p1f1"
would be valid for this argument

\item[\code{scenario}] character vector of future climate scenarios. Each value
in this vector must be either historical, ssp245, ssp370, or ssp585
\end{ldescription}
\end{Arguments}
%
\begin{Value}
list of length two - the tmax raster stack and the tmin raster
stack
\end{Value}
%
\begin{Examples}
\begin{ExampleCode}
## Not run: 
# Assuming you used the download_loca2() function and placed the output
# in a folder called LOCA2
loca_t_rast("LOCA2/ACCESS-CM2/0p0625deg/r1i1p1f1",
            c("historical", "ssp245"))

## End(Not run)
\end{ExampleCode}
\end{Examples}
